\chapter{Mantenimiento civil y militar}\label{ch:23}
\objetivos{Identificar el régimen por actividad y autoridad; comprender la armonización EMAR y evitar trasladar automáticamente aprobaciones civiles al sistema militar.}
\section{Primero, el ámbito}
Una aeronave utilizada en determinadas actividades estatales puede quedar fuera del ámbito ordinario del Reglamento 2018/1139 en las condiciones de su artículo 2.3(a). Este punto y las posibilidades del propio artículo se leen completos. La exclusión no significa ausencia de requisitos de seguridad; conduce a identificar el sistema y autoridad correspondientes.

El material docente militar explica fragmentación nacional, cooperación y armonización. Contiene una referencia al «artículo 1.2» del reglamento base para actividades militares que no se adopta aquí como localizador vigente: se utiliza el artículo 2 del texto consultado para estudiar alcance y exclusiones.
\section{Aeronavegabilidad de la Defensa en España}
El Real Decreto 866/2015 aprueba el Reglamento de Aeronavegabilidad de la Defensa (RAD). El texto consolidado consultado en BOE identifica el ámbito y autoridades de ese sistema, junto con certificados y procedimientos. No se utiliza una convocatoria de empleo público como si fuese este reglamento: el archivo BOE-A-2022-21487 aportado es una resolución de proceso selectivo.

Para un caso español se identifica primero si la aeronave y actividad están dentro del RAD, qué autoridad y organización actúan y qué documentos técnicos y aprobaciones son aplicables. La similitud de una denominación con Parte 145 civil no demuestra por sí sola reconocimiento jurídico recíproco.
\section{MAWA, EDA y EMAR}
La Agencia Europea de Defensa facilita, mediante el foro MAWA, armonización entre autoridades militares nacionales. Los EMAR cubren ámbitos como certificación inicial, mantenimiento, licencias, formación y gestión de aeronavegabilidad. Proponen un marco armonizado, pero la autoridad militar nacional conserva su función y su implementación depende de decisiones nacionales.

La explicación oficial de EDA señala que el foro no tiene potestad para imponer reglamentos a los Estados participantes. Un EMAR no se convierte automáticamente en un reglamento de la Unión directamente aplicable por su título. Se investiga su adopción o implementación en el sistema pertinente.
\begin{table}[htbp]\centering\small\caption{Comparación de preguntas, no equivalencia de aprobaciones.}\label{tab:militar}
\begin{tabular}{p{.26\linewidth}p{.3\linewidth}p{.31\linewidth}}\toprule
Aspecto & Contexto civil europeo & Contexto militar\\\midrule
Ámbito & 2018/1139 y reglas de desarrollo & Marco nacional y actividad concreta\\[4pt]
Certificación & Parte 21 y base aplicable & RAD/requisitos militares adoptados\\[4pt]
Mantenimiento & Partes de 1321/2014 según población & Aprobaciones y requisitos del sistema militar\\[4pt]
Armonización & Sistema de la Unión y acuerdos aplicables & MAWA/EMAR y reconocimiento nacional\\\bottomrule
\end{tabular}\end{table}
\section{Particularidades operativas}
Las misiones militares pueden introducir perfiles de carga, utilización, despliegue y soporte distintos de los civiles. Esto modifica necesidades técnicas y logísticas; no autoriza a ignorar el sistema de aeronavegabilidad. Las condiciones de misión deben trasladarse a espectros, programas y datos de vida de forma controlada.

Una misma familia de motor puede tener configuraciones y limitaciones distintas según instalación. La experiencia de una flota no se transfiere sin evaluar diferencias de carga, ambiente y mantenimiento. La confidencialidad puede limitar disponibilidad de datos, pero no se sustituye con valores inventados.
\section{Reconocimiento y cooperación}
El reconocimiento de autoridades o aprobaciones busca evitar duplicaciones cuando existe una base acordada. No es equivalente a asumir que cualquier certificado extranjero tiene validez automática. Se examinan alcance, productos, funciones y condiciones de reconocimiento.

En contratos multinacionales conviene separar autoridad del diseño, producción, operación y mantenimiento y documentar cómo se distribuyen responsabilidades. La claridad evita que una interfaz contractual o documental deje una tarea sin responsable.
\caso{M23. Aeronave similar, sistema distinto}{Una organización civil recibe una solicitud para intervenir una aeronave militar derivada de un modelo civil. Antes de aceptar, identifica autoridad, configuración, aprobación y datos militares aplicables y cualquier reconocimiento existente. La semejanza de piezas o la experiencia con el modelo civil es evidencia de capacidad potencial, no prueba de autorización para ejecutar y certificar esa intervención.}
\section{Armonización militar frente a legislación nacional}
EDA y el foro MAWA impulsaron requisitos EMAR para favorecer la armonización militar, pero ello no supone que una agencia supranacional certifique automáticamente todas las aeronaves de Defensa ni que EMAR sea directamente obligatorio sin incorporación nacional. Para un caso español deben comprobarse el ámbito y los textos militares nacionales aplicables, entre ellos el Reglamento de Aeronavegabilidad de la Defensa.

El archivo BOE-A-2022-21487 aportado en el paquete corresponde a una convocatoria de personal, no al reglamento militar. Siempre se verifica identificador, título, objeto, fecha y fuente oficial, en lugar de confiar en el nombre del PDF.


\section{Similitud técnica y régimen jurídico}
Un derivado militar puede compartir componentes con un tipo civil. La semejanza ayuda a comprender producto, pero no decide autoridad, aprobación ni reconocimiento. El artículo 2 del 2018/1139 delimita ámbito y exclusiones para actividades estatales concretas; no se convierte en exclusión general de toda aeronave pública de cualquier norma civil.

El RD 866/2015 aprueba el Reglamento de Aeronavegabilidad de la Defensa. Su uso exige identificar producto, actividad y disposiciones pertinentes. BOE-A-2022-21487 es una convocatoria de empleo público y no proporciona esa base.
\section{Armonización y reconocimiento}
MAWA es foro de armonización en el contexto de EDA. EMAR no se convierte por mera publicación en reglamento directamente obligatorio e idéntico para todos los Estados. Se comprueba implementación nacional y reconocimiento.

El reconocimiento tiene alcance y condiciones. Una aprobación civil o militar extranjera no es pase universal para cualquier tarea. El expediente identifica documento de aceptación, productos, funciones y límites. Si no consta, queda pendiente sin confundirlo con incapacidad técnica del taller.
\Needspace{.48\textheight}
\section{Configuración y misión}
Perfiles de misión, modificaciones, equipos y ambiente pueden diferir del derivado civil. Contadores y cargas responden a datos aplicables, sin heredar automáticamente todos los límites ni eliminar los que permanecen. Acceso a información también requiere control dentro del régimen.

Disciplina técnica y trazabilidad mantienen importancia en ambos ámbitos. Cambia el marco que asigna responsabilidades y reconoce evidencia; no existe un procedimiento único deducible para todos los productos militares.
\repaso{¿Puede un manual civil aportar información a un estudio militar?}{Puede, con alcance y configuración identificados. Su uso como dato aplicable y la certificación requieren el marco y aprobaciones correspondientes. La referencia por sí sola no crea reconocimiento.}

\resumen{La actividad y el ámbito determinan el sistema. El RAD es una referencia española concreta. EMAR busca armonización mediante implementación nacional; no crea por sí solo una autoridad supranacional equivalente a EASA. Reconocimiento y similitud técnica son conceptos distintos.}
\fuentes{2018/1139, art. 2 \cite{EU1139}; RD 866/2015 \cite{RAD}; EDA, MAWA Forum \cite{EDA}; presentación militar aportada, pp. 1--21.}
